% --- Archon physics formalization source begin ---
% archon:physics
% archon:covers PhyXMiniProblems/problem_phyx_mini_0514.lean
% archon:source-report reports/phyx_mini/problem_phyx_mini_0514.source.json

\chapter{Physics problem phyx\_mini\_0514}
\label{ch:PhyXMiniProblems_problem_phyx_mini_0514}

\paragraph{Problem source.}
\#\# Physical scenario

Many of the stars in the sky are actually binary stars, in which two stars orbit about their common center of mass. If the orbital speeds of the stars are high enough, the motion of the stars can be detected by the Doppler shifts of the light they emit. Stars for which this is the case are called \textbackslash{}emph\{spectroscopic binary stars\}. Figure shows the simplest case of a spectroscopic binary star: two identical stars, each with mass \textbackslash{}( m \textbackslash{}), orbiting their center of mass in a circle of radius \textbackslash{}( R \textbackslash{}). The plane of the stars' orbits is edge-on to the line of sight of an observer on the earth.

\#\# Question

The light from each star in the binary system varies from its maximum frequency to its minimum frequency and back again in 11.0 days. Determine the orbital radius \textbackslash{}( R \textbackslash{}) and the mass \textbackslash{}( m \textbackslash{}) of each star. Give your answer for \textbackslash{}( m \textbackslash{}) in kilograms and as a multiple of the mass of the sun, \textbackslash{}( 1.99 \textbackslash{}times 10\textasciicircum{}\{30\} \textbackslash{}) kg. Compare the value of \textbackslash{}( R \textbackslash{}) to the distance from the earth to the sun, \textbackslash{}( 1.50 \textbackslash{}times 10\textasciicircum{}\{11\} \textbackslash{}) m.

\#\# Answer choices

- A: A: \textbackslash{}( 4.56 \textbackslash{}times 10\textasciicircum{}4 \textbackslash{}text\{ m/s\} \textbackslash{})

- B: B: \textbackslash{}( 4.17 \textbackslash{}times 10\textasciicircum{}4 \textbackslash{}text\{ m/s\} \textbackslash{})

- C: C: \textbackslash{}( 3.62 \textbackslash{}times 10\textasciicircum{}4 \textbackslash{}text\{ m/s\} \textbackslash{})

- D: D: \textbackslash{}( 3.94 \textbackslash{}times 10\textasciicircum{}4 \textbackslash{}text\{ m/s\} \textbackslash{})

\#\# Figure caption (auxiliary; use the image as primary evidence)

The image shows a circular path with two objects, each labeled with "m," located on opposite sides of the circle. The circle is drawn in blue, and a central point is connected to the circle with a line labeled "R," indicating the radius. The objects are connected with green arrows pointing in opposite directions along the circle's path.

To the left of the circle, there is a large, leftward-pointing arrow labeled "To the earth," indicating a direction or force toward the Earth. The arrangement suggests a scenario involving circular motion with respect to the Earth.

\#\# Dataset metadata

Relativity; Physical Model Grounding Reasoning, Multi-Formula Reasoning

\paragraph{Recorded answer/context.}
D: D: \textbackslash{}( 3.94 \textbackslash{}times 10\textasciicircum{}4 \textbackslash{}text\{ m/s\} \textbackslash{})

\paragraph{Figure/image path.}
/root/proposal\_for\_physic/science-mango/phyx\_mini\_run/phyx\_data/test\_image/514.png

\paragraph{Formalization target.}
create a compiling Lean file with sorry bodies at `PhyXMiniProblems/problem\_phyx\_mini\_0514.lean`. The Lean declarations must preserve the physical quantities, dimensions or dimensional roles, figure labels, governing-law hypotheses, and final relation expressed by this problem.
Use Mathlib/Physlib names found through LeanExplore where available. If a physics API is missing, introduce faithful local abstractions rather than scalar placeholder aliases.

\begin{theorem}[Physics formalization target]
\label{thm:physics:phyx_mini_0514:target}
\lean{PhyXMiniProblems.ProblemPhyXMini0514.spectroscopicBinary_orbitalRadius_and_stellarMass}
\uses{def:physics:phyx-mini-0514:phyxminiproblems-problemphyxmini0514-lengthinmeters, def:physics:phyx-mini-0514:phyxminiproblems-problemphyxmini0514-timeinseconds, def:physics:phyx-mini-0514:phyxminiproblems-problemphyxmini0514-massinkilograms, def:physics:phyx-mini-0514:phyxminiproblems-problemphyxmini0514-speedinmeterspersecond, def:physics:phyx-mini-0514:phyxminiproblems-problemphyxmini0514-gravitationalconstantinsi, def:physics:phyx-mini-0514:phyxminiproblems-problemphyxmini0514-spectroscopicbinarysetup, def:physics:phyx-mini-0514:phyxminiproblems-problemphyxmini0514-matchesspectroscopicbinaryscenario, def:physics:phyx-mini-0514:phyxminiproblems-problemphyxmini0514-matchesproblemreadouts, def:physics:phyx-mini-0514:phyxminiproblems-problemphyxmini0514-matchesprimarybinarystarfigure, def:physics:phyx-mini-0514:phyxminiproblems-problemphyxmini0514-previousdopplerspeedresult, def:physics:phyx-mini-0514:phyxminiproblems-problemphyxmini0514-usestextbookgravitationalconstant, def:physics:phyx-mini-0514:phyxminiproblems-problemphyxmini0514-hasphysicalbinaryparameters, def:physics:phyx-mini-0514:phyxminiproblems-problemphyxmini0514-satisfiesspectroscopicdopplerlaws, def:physics:phyx-mini-0514:phyxminiproblems-problemphyxmini0514-satisfiesuniformcircularkinematics, def:physics:phyx-mini-0514:phyxminiproblems-problemphyxmini0514-satisfiesnewtonianbinarydynamics}
The assigned autoformalize agent should translate this physics problem into Lean declarations in the covered file, with theorem and lemma proof bodies written as `by sorry`.
\end{theorem}
\begin{proof}
This is an autoformalization task, not a proof task. Produce statements that can later be proved by the physics prover without weakening the physical model.
\end{proof}

% --- Archon declaration topology begin ---
\section{Declaration topology}
\begin{definition}[Length Quantity]
  \label{def:physics:phyx-mini-0514:phyxminiproblems-problemphyxmini0514-lengthquantity}
  \lean{PhyXMiniProblems.ProblemPhyXMini0514.LengthQuantity}
  A nonnegative physical length, independent of a choice of unit
\end{definition}
\begin{proof}
  This is the declared model definition of Length Quantity.
\end{proof}
\begin{definition}[Time Quantity]
  \label{def:physics:phyx-mini-0514:phyxminiproblems-problemphyxmini0514-timequantity}
  \lean{PhyXMiniProblems.ProblemPhyXMini0514.TimeQuantity}
  A nonnegative physical time interval
\end{definition}
\begin{proof}
  This is the declared model definition of Time Quantity.
\end{proof}
\begin{definition}[Mass Quantity]
  \label{def:physics:phyx-mini-0514:phyxminiproblems-problemphyxmini0514-massquantity}
  \lean{PhyXMiniProblems.ProblemPhyXMini0514.MassQuantity}
  A nonnegative physical mass
\end{definition}
\begin{proof}
  This is the declared model definition of Mass Quantity.
\end{proof}
\begin{definition}[Frequency Quantity]
  \label{def:physics:phyx-mini-0514:phyxminiproblems-problemphyxmini0514-frequencyquantity}
  \lean{PhyXMiniProblems.ProblemPhyXMini0514.FrequencyQuantity}
  A nonnegative physical frequency, with dimension inverse time
\end{definition}
\begin{proof}
  This is the declared model definition of Frequency Quantity.
\end{proof}
\begin{definition}[Force Quantity]
  \label{def:physics:phyx-mini-0514:phyxminiproblems-problemphyxmini0514-forcequantity}
  \lean{PhyXMiniProblems.ProblemPhyXMini0514.ForceQuantity}
  A nonnegative force, whose coherent SI unit is the newton
\end{definition}
\begin{proof}
  This is the declared model definition of Force Quantity.
\end{proof}
\begin{definition}[Gravitational Constant Quantity]
  \label{def:physics:phyx-mini-0514:phyxminiproblems-problemphyxmini0514-gravitationalconstantquantity}
  \lean{PhyXMiniProblems.ProblemPhyXMini0514.GravitationalConstantQuantity}
  The Newtonian gravitational constant, of physical dimension length³ * mass⁻¹ * time⁻²
\end{definition}
\begin{proof}
  This is the declared model definition of Gravitational Constant Quantity.
\end{proof}
\begin{definition}[length In Meters]
  \label{def:physics:phyx-mini-0514:phyxminiproblems-problemphyxmini0514-lengthinmeters}
  \lean{PhyXMiniProblems.ProblemPhyXMini0514.lengthInMeters}
  \uses{def:physics:phyx-mini-0514:phyxminiproblems-problemphyxmini0514-lengthquantity}
  Read a physical length in SI metres
\end{definition}
\begin{proof}
  This is the declared model definition of length In Meters.
\end{proof}
\begin{definition}[time In Seconds]
  \label{def:physics:phyx-mini-0514:phyxminiproblems-problemphyxmini0514-timeinseconds}
  \lean{PhyXMiniProblems.ProblemPhyXMini0514.timeInSeconds}
  \uses{def:physics:phyx-mini-0514:phyxminiproblems-problemphyxmini0514-timequantity}
  Read a physical time interval in SI seconds
\end{definition}
\begin{proof}
  This is the declared model definition of time In Seconds.
\end{proof}
\begin{definition}[mass In Kilograms]
  \label{def:physics:phyx-mini-0514:phyxminiproblems-problemphyxmini0514-massinkilograms}
  \lean{PhyXMiniProblems.ProblemPhyXMini0514.massInKilograms}
  \uses{def:physics:phyx-mini-0514:phyxminiproblems-problemphyxmini0514-massquantity}
  Read a physical mass in SI kilograms
\end{definition}
\begin{proof}
  This is the declared model definition of mass In Kilograms.
\end{proof}
\begin{definition}[speed In Meters Per Second]
  \label{def:physics:phyx-mini-0514:phyxminiproblems-problemphyxmini0514-speedinmeterspersecond}
  \lean{PhyXMiniProblems.ProblemPhyXMini0514.speedInMetersPerSecond}
  Read a nonnegative physical speed in SI metres per second
\end{definition}
\begin{proof}
  This is the declared model definition of speed In Meters Per Second.
\end{proof}
\begin{definition}[frequency In Hertz]
  \label{def:physics:phyx-mini-0514:phyxminiproblems-problemphyxmini0514-frequencyinhertz}
  \lean{PhyXMiniProblems.ProblemPhyXMini0514.frequencyInHertz}
  \uses{def:physics:phyx-mini-0514:phyxminiproblems-problemphyxmini0514-frequencyquantity}
  Read a physical frequency in SI hertz
\end{definition}
\begin{proof}
  This is the declared model definition of frequency In Hertz.
\end{proof}
\begin{definition}[force In Newtons]
  \label{def:physics:phyx-mini-0514:phyxminiproblems-problemphyxmini0514-forceinnewtons}
  \lean{PhyXMiniProblems.ProblemPhyXMini0514.forceInNewtons}
  \uses{def:physics:phyx-mini-0514:phyxminiproblems-problemphyxmini0514-forcequantity}
  Read a physical force in SI newtons
\end{definition}
\begin{proof}
  This is the declared model definition of force In Newtons.
\end{proof}
\begin{definition}[gravitational Constant In SI]
  \label{def:physics:phyx-mini-0514:phyxminiproblems-problemphyxmini0514-gravitationalconstantinsi}
  \lean{PhyXMiniProblems.ProblemPhyXMini0514.gravitationalConstantInSI}
  \uses{def:physics:phyx-mini-0514:phyxminiproblems-problemphyxmini0514-gravitationalconstantquantity}
  Read the gravitational constant in m³ kg⁻¹ s⁻²
\end{definition}
\begin{proof}
  This is the declared model definition of gravitational Constant In SI.
\end{proof}
\begin{definition}[speed Of Light In Meters Per Second]
  \label{def:physics:phyx-mini-0514:phyxminiproblems-problemphyxmini0514-speedoflightinmeterspersecond}
  \lean{PhyXMiniProblems.ProblemPhyXMini0514.speedOfLightInMetersPerSecond}
  Physlib's exact vacuum speed of light, read in metres per second
\end{definition}
\begin{proof}
  This is the declared model definition of speed Of Light In Meters Per Second.
\end{proof}
\begin{definition}[Binary Star]
  \label{def:physics:phyx-mini-0514:phyxminiproblems-problemphyxmini0514-binarystar}
  \lean{PhyXMiniProblems.ProblemPhyXMini0514.BinaryStar}
  The two otherwise identical stars of the binary system
\end{definition}
\begin{proof}
  This is the declared model definition of Binary Star.
\end{proof}
\begin{definition}[other Star]
  \label{def:physics:phyx-mini-0514:phyxminiproblems-problemphyxmini0514-otherstar}
  \lean{PhyXMiniProblems.ProblemPhyXMini0514.otherStar}
  \uses{def:physics:phyx-mini-0514:phyxminiproblems-problemphyxmini0514-binarystar}
  The companion of a selected star
\end{definition}
\begin{proof}
  This is the declared model definition of other Star.
\end{proof}
\begin{definition}[Orbit Shape]
  \label{def:physics:phyx-mini-0514:phyxminiproblems-problemphyxmini0514-orbitshape}
  \lean{PhyXMiniProblems.ProblemPhyXMini0514.OrbitShape}
  Shape of the stellar trajectories about the common center
\end{definition}
\begin{proof}
  This is the declared model definition of Orbit Shape.
\end{proof}
\begin{definition}[Viewing Orientation]
  \label{def:physics:phyx-mini-0514:phyxminiproblems-problemphyxmini0514-viewingorientation}
  \lean{PhyXMiniProblems.ProblemPhyXMini0514.ViewingOrientation}
  Orientation of the orbital plane relative to the terrestrial observer
\end{definition}
\begin{proof}
  This is the declared model definition of Viewing Orientation.
\end{proof}
\begin{definition}[Figure Position]
  \label{def:physics:phyx-mini-0514:phyxminiproblems-problemphyxmini0514-figureposition}
  \lean{PhyXMiniProblems.ProblemPhyXMini0514.FigurePosition}
  Location of a star in the particular snapshot shown by the bitmap
\end{definition}
\begin{proof}
  This is the declared model definition of Figure Position.
\end{proof}
\begin{definition}[Horizontal Direction]
  \label{def:physics:phyx-mini-0514:phyxminiproblems-problemphyxmini0514-horizontaldirection}
  \lean{PhyXMiniProblems.ProblemPhyXMini0514.HorizontalDirection}
  Horizontal direction used by the Earth arrow and tangential arrows
\end{definition}
\begin{proof}
  This is the declared model definition of Horizontal Direction.
\end{proof}
\begin{definition}[Binary Star Figure]
  \label{def:physics:phyx-mini-0514:phyxminiproblems-problemphyxmini0514-binarystarfigure}
  \lean{PhyXMiniProblems.ProblemPhyXMini0514.BinaryStarFigure}
  \uses{def:physics:phyx-mini-0514:phyxminiproblems-problemphyxmini0514-binarystar, def:physics:phyx-mini-0514:phyxminiproblems-problemphyxmini0514-figureposition, def:physics:phyx-mini-0514:phyxminiproblems-problemphyxmini0514-horizontaldirection}
  Qualitative information transcribed from the primary bitmap. The radius segment is a geometrical radius of the blue orbital circle; it is not a local definition of the requested numerical radius
\end{definition}
\begin{proof}
  This is the declared model definition of Binary Star Figure.
\end{proof}
\begin{definition}[Spectroscopic Binary Setup]
  \label{def:physics:phyx-mini-0514:phyxminiproblems-problemphyxmini0514-spectroscopicbinarysetup}
  \lean{PhyXMiniProblems.ProblemPhyXMini0514.SpectroscopicBinarySetup}
  \uses{def:physics:phyx-mini-0514:phyxminiproblems-problemphyxmini0514-lengthquantity, def:physics:phyx-mini-0514:phyxminiproblems-problemphyxmini0514-timequantity, def:physics:phyx-mini-0514:phyxminiproblems-problemphyxmini0514-massquantity, def:physics:phyx-mini-0514:phyxminiproblems-problemphyxmini0514-frequencyquantity, def:physics:phyx-mini-0514:phyxminiproblems-problemphyxmini0514-forcequantity, def:physics:phyx-mini-0514:phyxminiproblems-problemphyxmini0514-gravitationalconstantquantity, def:physics:phyx-mini-0514:phyxminiproblems-problemphyxmini0514-binarystar, def:physics:phyx-mini-0514:phyxminiproblems-problemphyxmini0514-orbitshape, def:physics:phyx-mini-0514:phyxminiproblems-problemphyxmini0514-viewingorientation, def:physics:phyx-mini-0514:phyxminiproblems-problemphyxmini0514-binarystarfigure}
  Physical quantities belonging to the binary system and its observation. None of the radius or mass fields is defined from a requested numerical answer; their relations are supplied only by premise interfaces below
\end{definition}
\begin{proof}
  This is the declared model definition of Spectroscopic Binary Setup.
\end{proof}
\begin{definition}[Matches Spectroscopic Binary Scenario]
  \label{def:physics:phyx-mini-0514:phyxminiproblems-problemphyxmini0514-matchesspectroscopicbinaryscenario}
  \lean{PhyXMiniProblems.ProblemPhyXMini0514.MatchesSpectroscopicBinaryScenario}
  \uses{def:physics:phyx-mini-0514:phyxminiproblems-problemphyxmini0514-lengthinmeters, def:physics:phyx-mini-0514:phyxminiproblems-problemphyxmini0514-otherstar, def:physics:phyx-mini-0514:phyxminiproblems-problemphyxmini0514-spectroscopicbinarysetup}
  The equal-star, circular, edge-on physical scenario stated in the prose
\end{definition}
\begin{proof}
  This is the declared model definition of Matches Spectroscopic Binary Scenario.
\end{proof}
\begin{definition}[Matches Problem Readouts]
  \label{def:physics:phyx-mini-0514:phyxminiproblems-problemphyxmini0514-matchesproblemreadouts}
  \lean{PhyXMiniProblems.ProblemPhyXMini0514.MatchesProblemReadouts}
  \uses{def:physics:phyx-mini-0514:phyxminiproblems-problemphyxmini0514-lengthinmeters, def:physics:phyx-mini-0514:phyxminiproblems-problemphyxmini0514-timeinseconds, def:physics:phyx-mini-0514:phyxminiproblems-problemphyxmini0514-massinkilograms, def:physics:phyx-mini-0514:phyxminiproblems-problemphyxmini0514-frequencyinhertz, def:physics:phyx-mini-0514:phyxminiproblems-problemphyxmini0514-spectroscopicbinarysetup}
  The measured frequency pattern and the comparison constants printed in the question. A maximum-to-minimum-to-maximum cycle is one complete orbit in the edge-on circular model. No numerical radius or stellar mass is recorded
\end{definition}
\begin{proof}
  This is the declared model definition of Matches Problem Readouts.
\end{proof}
\begin{definition}[Matches Primary Binary Star Figure]
  \label{def:physics:phyx-mini-0514:phyxminiproblems-problemphyxmini0514-matchesprimarybinarystarfigure}
  \lean{PhyXMiniProblems.ProblemPhyXMini0514.MatchesPrimaryBinaryStarFigure}
  \uses{def:physics:phyx-mini-0514:phyxminiproblems-problemphyxmini0514-spectroscopicbinarysetup}
  Primary-image evidence, kept separate from dynamical laws
\end{definition}
\begin{proof}
  This is the declared model definition of Matches Primary Binary Star Figure.
\end{proof}
\begin{definition}[Answer Choice]
  \label{def:physics:phyx-mini-0514:phyxminiproblems-problemphyxmini0514-answerchoice}
  \lean{PhyXMiniProblems.ProblemPhyXMini0514.AnswerChoice}
  Labels of the four speed choices supplied by the inconsistent source
\end{definition}
\begin{proof}
  This is the declared model definition of Answer Choice.
\end{proof}
\begin{definition}[answer Speed In Meters Per Second]
  \label{def:physics:phyx-mini-0514:phyxminiproblems-problemphyxmini0514-answerspeedinmeterspersecond}
  \lean{PhyXMiniProblems.ProblemPhyXMini0514.answerSpeedInMetersPerSecond}
  \uses{def:physics:phyx-mini-0514:phyxminiproblems-problemphyxmini0514-answerchoice}
  Metres-per-second readout printed beside each answer label
\end{definition}
\begin{proof}
  This is the declared model definition of answer Speed In Meters Per Second.
\end{proof}
\begin{definition}[Previous Doppler Speed Result]
  \label{def:physics:phyx-mini-0514:phyxminiproblems-problemphyxmini0514-previousdopplerspeedresult}
  \lean{PhyXMiniProblems.ProblemPhyXMini0514.PreviousDopplerSpeedResult}
  \uses{def:physics:phyx-mini-0514:phyxminiproblems-problemphyxmini0514-speedinmeterspersecond, def:physics:phyx-mini-0514:phyxminiproblems-problemphyxmini0514-spectroscopicbinarysetup, def:physics:phyx-mini-0514:phyxminiproblems-problemphyxmini0514-answerspeedinmeterspersecond}
  Previous spectroscopic analysis supplying the speed recorded by answer D. This is explicitly a prior calibration because the current prose asks for R and m, and the source omits the numerical frequency extrema needed to derive the speed afresh
\end{definition}
\begin{proof}
  This is the declared model definition of Previous Doppler Speed Result.
\end{proof}
\begin{definition}[Uses Textbook Gravitational Constant]
  \label{def:physics:phyx-mini-0514:phyxminiproblems-problemphyxmini0514-usestextbookgravitationalconstant}
  \lean{PhyXMiniProblems.ProblemPhyXMini0514.UsesTextbookGravitationalConstant}
  \uses{def:physics:phyx-mini-0514:phyxminiproblems-problemphyxmini0514-gravitationalconstantinsi, def:physics:phyx-mini-0514:phyxminiproblems-problemphyxmini0514-spectroscopicbinarysetup}
  Standard textbook value G = 6.67 × 10⁻¹¹ m³ kg⁻¹ s⁻²
\end{definition}
\begin{proof}
  This is the declared model definition of Uses Textbook Gravitational Constant.
\end{proof}
\begin{definition}[Has Physical Binary Parameters]
  \label{def:physics:phyx-mini-0514:phyxminiproblems-problemphyxmini0514-hasphysicalbinaryparameters}
  \lean{PhyXMiniProblems.ProblemPhyXMini0514.HasPhysicalBinaryParameters}
  \uses{def:physics:phyx-mini-0514:phyxminiproblems-problemphyxmini0514-lengthinmeters, def:physics:phyx-mini-0514:phyxminiproblems-problemphyxmini0514-timeinseconds, def:physics:phyx-mini-0514:phyxminiproblems-problemphyxmini0514-massinkilograms, def:physics:phyx-mini-0514:phyxminiproblems-problemphyxmini0514-speedinmeterspersecond, def:physics:phyx-mini-0514:phyxminiproblems-problemphyxmini0514-frequencyinhertz, def:physics:phyx-mini-0514:phyxminiproblems-problemphyxmini0514-gravitationalconstantinsi, def:physics:phyx-mini-0514:phyxminiproblems-problemphyxmini0514-speedoflightinmeterspersecond, def:physics:phyx-mini-0514:phyxminiproblems-problemphyxmini0514-spectroscopicbinarysetup}
  Positivity and subluminality conditions for the physical model
\end{definition}
\begin{proof}
  This is the declared model definition of Has Physical Binary Parameters.
\end{proof}
\begin{definition}[Satisfies Spectroscopic Doppler Laws]
  \label{def:physics:phyx-mini-0514:phyxminiproblems-problemphyxmini0514-satisfiesspectroscopicdopplerlaws}
  \lean{PhyXMiniProblems.ProblemPhyXMini0514.SatisfiesSpectroscopicDopplerLaws}
  \uses{def:physics:phyx-mini-0514:phyxminiproblems-problemphyxmini0514-speedinmeterspersecond, def:physics:phyx-mini-0514:phyxminiproblems-problemphyxmini0514-frequencyinhertz, def:physics:phyx-mini-0514:phyxminiproblems-problemphyxmini0514-speedoflightinmeterspersecond, def:physics:phyx-mini-0514:phyxminiproblems-problemphyxmini0514-spectroscopicbinarysetup}
  Longitudinal relativistic Doppler shift for approaching and receding extrema. For an edge-on circular orbit, the line-of-sight speed amplitude equals the orbital speed. These laws contain no radius or stellar-mass answer
\end{definition}
\begin{proof}
  This is the declared model definition of Satisfies Spectroscopic Doppler Laws.
\end{proof}
\begin{definition}[Satisfies Uniform Circular Kinematics]
  \label{def:physics:phyx-mini-0514:phyxminiproblems-problemphyxmini0514-satisfiesuniformcircularkinematics}
  \lean{PhyXMiniProblems.ProblemPhyXMini0514.SatisfiesUniformCircularKinematics}
  \uses{def:physics:phyx-mini-0514:phyxminiproblems-problemphyxmini0514-lengthinmeters, def:physics:phyx-mini-0514:phyxminiproblems-problemphyxmini0514-timeinseconds, def:physics:phyx-mini-0514:phyxminiproblems-problemphyxmini0514-speedinmeterspersecond, def:physics:phyx-mini-0514:phyxminiproblems-problemphyxmini0514-spectroscopicbinarysetup}
  Uniform circular kinematics: in one orbital period, each star travels one circumference of its own radius about the common center of mass
\end{definition}
\begin{proof}
  This is the declared model definition of Satisfies Uniform Circular Kinematics.
\end{proof}
\begin{definition}[Satisfies Newtonian Binary Dynamics]
  \label{def:physics:phyx-mini-0514:phyxminiproblems-problemphyxmini0514-satisfiesnewtonianbinarydynamics}
  \lean{PhyXMiniProblems.ProblemPhyXMini0514.SatisfiesNewtonianBinaryDynamics}
  \uses{def:physics:phyx-mini-0514:phyxminiproblems-problemphyxmini0514-lengthinmeters, def:physics:phyx-mini-0514:phyxminiproblems-problemphyxmini0514-massinkilograms, def:physics:phyx-mini-0514:phyxminiproblems-problemphyxmini0514-speedinmeterspersecond, def:physics:phyx-mini-0514:phyxminiproblems-problemphyxmini0514-forceinnewtons, def:physics:phyx-mini-0514:phyxminiproblems-problemphyxmini0514-gravitationalconstantinsi, def:physics:phyx-mini-0514:phyxminiproblems-problemphyxmini0514-otherstar, def:physics:phyx-mini-0514:phyxminiproblems-problemphyxmini0514-spectroscopicbinarysetup}
  Newtonian two-body dynamics in the center-of-mass frame. Antipodal geometry makes the center-to-center separation the sum of the two orbital radii; inverse-square gravity supplies the mutual force, and that same force is the centripetal force for each circular trajectory. These are unsolved laws, not the requested formulas for R or m
\end{definition}
\begin{proof}
  This is the declared model definition of Satisfies Newtonian Binary Dynamics.
\end{proof}
\begin{lemma}[orbit Radius eq speed mul period div two pi]
  \label{lem:physics:phyx-mini-0514:phyxminiproblems-problemphyxmini0514-orbitradius-eq-speed-mul-period-div-two-pi}
  \lean{PhyXMiniProblems.ProblemPhyXMini0514.orbitRadius_eq_speed_mul_period_div_two_pi}
  \uses{def:physics:phyx-mini-0514:phyxminiproblems-problemphyxmini0514-lengthinmeters, def:physics:phyx-mini-0514:phyxminiproblems-problemphyxmini0514-timeinseconds, def:physics:phyx-mini-0514:phyxminiproblems-problemphyxmini0514-speedinmeterspersecond, def:physics:phyx-mini-0514:phyxminiproblems-problemphyxmini0514-spectroscopicbinarysetup, def:physics:phyx-mini-0514:phyxminiproblems-problemphyxmini0514-hasphysicalbinaryparameters, def:physics:phyx-mini-0514:phyxminiproblems-problemphyxmini0514-satisfiesuniformcircularkinematics}
  Circular kinematics determines each orbital radius from speed and period
\end{lemma}
\begin{proof}
  The conclusion follows from the governing relations and figure readouts named in the statement of orbit Radius eq speed mul period div two pi.
\end{proof}
\begin{lemma}[star Mass eq four radius speed sq div gravitational Constant]
  \label{lem:physics:phyx-mini-0514:phyxminiproblems-problemphyxmini0514-starmass-eq-four-radius-speed-sq-div-gravitationalconstant}
  \lean{PhyXMiniProblems.ProblemPhyXMini0514.starMass_eq_four_radius_speed_sq_div_gravitationalConstant}
  \uses{def:physics:phyx-mini-0514:phyxminiproblems-problemphyxmini0514-lengthinmeters, def:physics:phyx-mini-0514:phyxminiproblems-problemphyxmini0514-massinkilograms, def:physics:phyx-mini-0514:phyxminiproblems-problemphyxmini0514-speedinmeterspersecond, def:physics:phyx-mini-0514:phyxminiproblems-problemphyxmini0514-gravitationalconstantinsi, def:physics:phyx-mini-0514:phyxminiproblems-problemphyxmini0514-spectroscopicbinarysetup, def:physics:phyx-mini-0514:phyxminiproblems-problemphyxmini0514-matchesspectroscopicbinaryscenario, def:physics:phyx-mini-0514:phyxminiproblems-problemphyxmini0514-hasphysicalbinaryparameters, def:physics:phyx-mini-0514:phyxminiproblems-problemphyxmini0514-satisfiesnewtonianbinarydynamics}
  Equal masses and radii reduce Newtonian gravity plus centripetal motion to m = 4 R v² / G for either star
\end{lemma}
\begin{proof}
  The conclusion follows from the governing relations and figure readouts named in the statement of star Mass eq four radius speed sq div gravitational Constant.
\end{proof}
\begin{definition}[Recorded answer choice]
  \label{def:physics:phyx-mini-0514:phyxminiproblems-problemphyxmini0514-recordedanswerchoice}
  \lean{PhyXMiniProblems.ProblemPhyXMini0514.recordedAnswerChoice}
  \uses{def:physics:phyx-mini-0514:phyxminiproblems-problemphyxmini0514-answerchoice}
  The private dataset metadata records choice D; it is not a spectroscopic readout or a premise of the symbolic binary-star result.
\end{definition}
\begin{proof}
  This is the literal private metadata assignment; the block is retained because the Lean-aware graph scans this declaration.
\end{proof}
% --- Archon declaration topology end ---
% --- Archon physics formalization source end ---
